\chapter{Introduction}
%Maybe give a brief overview on Carbon Emissions 
% and the important of reducing them - Maybe too easy?
Reduction in the usage of energy in buildings is a crucial aspect 
if one wants to go towards a more sustainable future. 
It has long been known that buildings are responseible for X\% of the total energy 
consumption in the world, and that a large part of this energy 
is for HVAC systems [\textcolor{red}{MISSING REFERENCE}]
HVAC systems energy consumption has been directly related to the occupant behaviours of the 
buildings. This imposes a challenge, as one has to be able to understand the behaviour of the occupants 
in the buildings before we can gain the largest optimization in energy consumption. 
(\textcolor{red}{MISSING REFERENCE}) has shown that there is a large variability in the behavour of occupants such as 
thermal comfort, lighting \& ventiliation needs and precense hours. 
\bigskip

%Maybe write in the drawbacks on the current methods of heating
%i.e. assuming constant thermal comfort and presence
Within the context of energy consumption in buildings, heating has been estimated
to contribute to X\% [\textcolor{red}{MISSING REFERENCE}] of the consumption. 
An aspect of heating of a building is the byproduct of the thermal comfort of the occupants in the building, 
and therefore, if we can understand the thermal comforts of the occupants in the buildings, 
we can optimize the energy consumption drasitically. However, the thermal comfort curves
is a personal preference of the occupants, and we should therefore expect this to differer dependent 
on the individual buildings and the number of occupants. 

%Write out litterature already done on the topic and the drawbacks of these methods.
%Maybe Write out the benefits of using a HMM
\bigskip
This projects aims to understand when the occupants of a single offices are present,
and how much thermal comfort feedback does we need to collect from the occupants in order to be able to 
infer the true thermal comfort of the occupants in the building. 
The goal is to also build a codebase others and easyly download and extent.\\
We will do this by using a Hidden Markov Model (HMM) to model the state of the room 
%Write out why using CO2 data is a good proxy for the presence of the occupants in the building.


\bigskip
\textcolor{red}{WRITTEN OUTLINE OF EACH CHAPTER AND WHAT QUESTIONS THEY ANSWER}